% =====================================================================
% Appendix A.5 — Reassembly replicates
%
% Regenerated by: audit/scripts/appendix_A4_A5_tables.py  (exit 0)
% Data: four_floor/characterisation/rebuild1 .. rebuild5, five taps each
% Audit item: new in the appendix pass
%
% *** AUTHOR: READ THE ROUNDING NOTE BELOW BEFORE SUBMITTING ***
% The published floors 0.30 / 0.46 / 0.32 are 2x the 1-sigma values ROUNDED
% to two decimals (0.15, 0.23, 0.16). Doubling before rounding gives
% 0.31 / 0.45 / 0.31. This section states that convention explicitly so the
% appendix and the body do not disagree silently. If you would rather change
% the body's floors instead, every derived ratio in Chapter 4 moves with them
% (e.g. Floor 1 trace 14.2x becomes 13.8x). See the audit report.
%
% PACKAGES REQUIRED: booktabs, amsmath
% Labels prefixed app:reassembly / tab:reassembly.
% =====================================================================

\section{Reassembly replicates}
\label{app:reassembly}

The reassembly floors of $0.30$, $0.46$ and $0.32\%$ underpin every detection
claim in Chapter~4: a frequency shift counts as damage only when it exceeds the
floor for that mode. Table~4.3 gives the floors alone. The five teardown and
rebuild cycles they come from are given here.

Each cycle was a complete teardown and reassembly of the undamaged frame,
followed by five taps. The frame was returned to nominal condition each time, so
the spread across cycles measures how much the modal frequencies move under
reassembly alone, with no damage present.

\begin{table}[htbp]
\centering
\caption{Modal frequencies after each teardown and rebuild cycle. Each entry is
the mean over the five taps of that cycle.}
\label{tab:reassembly-cycles}
\begin{tabular}{lrrr}
\toprule
Cycle & $f_1$ (Hz) & $f_2$ (Hz) & $f_3$ (Hz) \\
\midrule
Cycle 1 & 2.9604 & 8.1468 & 12.1752 \\
Cycle 2 & 2.9585 & 8.1440 & 12.1867 \\
Cycle 3 & 2.9557 & 8.1280 & 12.1927 \\
Cycle 4 & 2.9638 & 8.1768 & 12.1552 \\
Cycle 5 & 2.9675 & 8.1610 & 12.2059 \\
\bottomrule
\end{tabular}
\end{table}

\subsection*{Deriving the floors}

The floor for each mode is twice the standard deviation of that mode across the
five cycles, expressed as a percentage of its across-cycle mean.

\begin{table}[htbp]
\centering
\caption{Derivation of the reassembly floors. The $1\sigma$ column is the
across-cycle standard deviation as a percentage of the across-cycle mean.}
\label{tab:reassembly-floors}
\begin{tabular}{lrrrrr}
\toprule
Mode & Cycles & Mean (Hz) & SD (Hz) & $1\sigma$ (\%) & $2\sigma$ (\%) \\
\midrule
$f_1$ & 5 & 2.9612 & 0.00458 & 0.155 & 0.31 \\
$f_2$ & 5 & 8.1513 & 0.01846 & 0.226 & 0.45 \\
$f_3$ & 5 & 12.1831 & 0.01914 & 0.157 & 0.31 \\
\bottomrule
\end{tabular}
\end{table}

\subsection*{Rounding convention}

The floors quoted throughout Chapter~4, $0.30$, $0.46$ and $0.32\%$, are twice
the $1\sigma$ values rounded to two decimals, that is $2 \times 0.15$,
$2 \times 0.23$ and $2 \times 0.16$. Doubling the unrounded values instead gives
$0.31$, $0.45$ and $0.31\%$, as in Table~\ref{tab:reassembly-floors}. The
difference is in the second decimal place and changes no detection outcome: the
smallest margin in the campaign, Floor~2 at trace grade, clears the first-mode
floor by $3.2$ times under the quoted convention and $3.1$ times under the other.
The quoted values are retained for consistency with Table~4.3.

\subsection*{What $n = 5$ supports}

Five cycles give four degrees of freedom, so these floors are operational
thresholds rather than formal statistical limits. A $2\sigma$ interval estimated
from five observations carries wide uncertainty in its own right, and no
distributional claim is made about the reassembly process. The floors are used in
Chapter~4 as a measured yardstick for what a rebuild moves, which is what a
deployed node would face, and not as a confidence bound. Where a margin is close
to its floor, the conclusion is stated as provisional for that reason.
